\section{Introduction}\label{sec:introduction}

A Hausdorff dimension identifies a critical power of distance, but that
power need not assign nonzero measure to a random curve. For chordal
Schramm--Loewner evolution $\SLE_\kappa$, $0<\kappa<8$, Beffara
\cite{Beffara2008} proved that the almost-sure dimension is
\[
 d=1+\frac\kappa8.
\]
Schramm asked whether the trace has a sigma-finite Hausdorff measure
and suggested $r^d\log\log(1/r)$ as a possible gauge
\cite[Problem~7.1 and the following paragraph]{Schramm2006}.
Rezaei subsequently proved that the critical power $r^d$ itself gives
zero Hausdorff measure \cite[Theorem~1.1]{Rezaei2018}.
The question therefore asks for a deterministic refinement of dimension
that measures the trace through arbitrary spatial covers.

We answer the existence question positively by constructing the scales
of the gauge from moments of a local trace mass. The gauge is explicit
in terms of deterministic integrals, although its asymptotics need not
have a regular form. The main mechanism uses finite moments of every
order in two ways: moderately large masses provide repeated opportunities
for an economical cover, while exceedances of a larger threshold have a
summable total cost. A sharp tail asymptotic is not needed.

\subsection{The result}

A \emph{Hausdorff gauge} is a continuous nondecreasing function
$h:[0,\infty)\to[0,\infty)$ with $h(0)=0$ and $h(r)>0$ for $r>0$.
We use diameters in the definition
\[
 \Haus^h(E)=\lim_{\delta\downarrow0}\Haus^h_\delta(E),\qquad
 \Haus^h_\delta(E)=\inf\left\{\sum_j h(\diam U_j):
 E\subset\bigcup_jU_j,\quad\diam U_j\le\delta\right\},
\]
where covers are countable and their sets may be arbitrary.

\begin{theorem}[An exact moment gauge]\label{thm:main}
Fix $0<\kappa<8$, and let $\gamma$ be chordal $\SLE_\kappa$ from
$0$ to $\infty$ in the upper half-plane $\HH$, parametrized by half-plane capacity $2t$.
There is a deterministic Hausdorff gauge $h_\kappa$, depending only on
$\kappa$, such that almost surely, simultaneously for every real
$0<s<t<\infty$,
\[
 0<\Haus^{h_\kappa}(\gamma([s,t]))<\infty.
\]
\end{theorem}

The gauge is the full infimum in \eqref{gauge:definition}, with
coefficients specified by the iterated integrals in \eqref{coef:formula}.
The exceptional event may depend on the fixed $\kappa$; the gauge
is independent of both the sample and the interval.
In particular, the trace outside its initial point has a nonzero
sigma-finite Hausdorff measure.

The upper argument measures the entire trace in a bounded region.
For $\Gamma=\gamma([0,\infty))$ and
$B_m=[-m,m]+i[0,m]$, Proposition~\ref{hd:upper} proves
\begin{equation}
 \EE\,\Haus^{h_\kappa}(\Gamma\cap B_m)<\infty
 \qquad(m\in\mathbb N).
 \label{intro:box-mean}
\end{equation}
Thus all these bounded-box measures are finite on one event, including
at real-boundary trace points. Positivity in Theorem~\ref{thm:main}
is asserted for the stated positive-time intervals.

\subsection{Multipoint estimates and one canonical mass}

The gauge construction needs a trace mass with moments of every order
and a way to compare its future in a small region with one fixed disk
model. We establish these inputs through a multipoint probability bound
and a conformally restarting mass. Both are also useful independently.
We first give a direct product-supermartingale proof of the fixed-order
estimate needed in an arbitrary physical domain.
For a conformal map $J:\HH\to D$, let $R_D(x)$ denote full conformal
radius: $R_D(x)=|F'(0)|$ when $F:\DD\to D$ is conformal,
$\DD$ is the unit disk, and $F(0)=x$.
Write $H_t$ for the surviving half-plane domain after the trace up to
time $t$, and $D_t=J(H_t)$. A point is \emph{alive} at $t$ when it
belongs to $D_t$; its evolving radius is $R_{D_t}(x)$.
Set $\alpha=2-d$ and $a=2/\kappa$.
For distinct interior points $x_1,\ldots,x_m$, choose caps satisfying
\[
 0<b_i\le R_D(x_i),\qquad b_i\le |x_i-x_j|\quad(j\ne i).
\]
Proposition~\ref{mp:estimate} proves that the probability that each point
reaches its own radius target $e_i>0$ while alive is at most
\begin{equation}
 C_{m,a}\prod_{i=1}^m\min\{1,(e_i/b_i)^\alpha\}.
 \label{intro:multipoint}
\end{equation}
The constant depends only on the fixed order and parameter. It is
uniform in the domain, the initial angles, and the order and times of
the hits. A completed point may be swallowed before another target is
reached. No sharp growth estimate in $m$ is claimed. Integrating the
collision singularities in this bound gives all the compact moments
used in the mass construction.

The second ingredient is a single adapted increasing measure process
$\mu_t$ on the open domain, starting at zero. In the half-plane,
Proposition~\ref{mass:construction} constructs it with nonatomic terminal
measure, local total-variation continuity, and increments carried by the
corresponding trace segments. Every nontrivial time interval has positive
increment mass by Proposition~\ref{mass:positive}.
To describe its conditional mean, let $g_S$ be the Loewner map of the
surviving half-plane domain $H_S$, with driving point $U_S$ in the
capacity normalization. For $x\in H_S$ put
$\theta_S(x)=\arg(g_S(x)-U_S)$ and
$G_S(x)=R_{H_S}(x)^{-\alpha}\sin^{8/\kappa-1}\theta_S(x)$;
put $G_S(x)=0$ outside $H_S$. Its remaining conditional mean is
\[
 \EE[\mu_\infty-\mu_S\mid\mathcal F_S]=G_S\,\dd A
\]
for every fixed stopping time $S$, including $S=\infty$ with
$G_\infty=0$. Here $\dd A$ is planar area. The equality holds as a
conditional measure kernel, so it includes nonnegative tests known
at $S$.

Conformal covariance weights the measure by $|J'|^d$.
At each fixed finite stopping time, Proposition~\ref{mass:restart}
identifies the actual future increment with the conformally transported
canonical mass of the fresh driver, simultaneously at every later time.
This permits comparison of the actual mass near a marked point with
one disk model after stopping at a prescribed conformal radius.
The lower bound also uses first-visit accounting: mass already
accumulated on the trace is not charged again at a later geometric return.

\subsection{Why the moments determine the gauge}

To bias the curve toward a marked interior point, stop when that point's
conformal radius reaches a prescribed smaller value and reweight the
stopped law by the ratio of its terminal and initial Green weights.
Paths that fail to reach the radius level receive weight zero.
Section~\ref{op:section} proves that the weight has mean one and that
these finite stopped laws are consistent as the radius decreases.
We call them the \emph{centered laws}.

Place the marked point at the origin of a unit disk and measure mass in
its radius-$1/4$ subdisk. The \emph{inner time} is the logarithmic radius
clock $\ell=(2a)^{-1}\log(R_0/R)$, where $R$ is the marked point's current
conformal radius. Let $L_\ell$ be the mass accumulated in that subdisk
by inner time $\ell$. Under the centered law, the marked point's angle
in the corresponding half-plane coordinates evolves as a diffusion
with invariant density proportional to
$\sin^{4a}\theta$. We average the initial angle against this distribution
and write
$\PP_{\mathrm{mix}}^*,\EE_{\mathrm{mix}}^*$ for this averaged centered
law and expectation. Then
\[
 L_\ell\uparrow L,\qquad 0<\EE_{\mathrm{mix}}^*L^n<\infty\quad(n\ge1).
\]
Section~\ref{coef:section} specifies these numbers by Wiener, angle and
area integrals, and then by product-normal coordinates and ordinary
differential equation limits. This deterministic infinite-limit
prescription fixes the normalization
of the moments without referring to an unspecified random measure.

For each order $n$, retain more than three quarters of its moment at a
finite level $\ell_n$. Repeat dyadic thresholds $A_k=2^k$ in a finite
batch with multiplicities $N_{nk}$ determined by the $n$th and
$(n+1)$st moments. Lemma~\ref{gauge:multiplicity} gives the complementary
bounds
\[
 \sum_kN_{nk}\PP_{\mathrm{mix}}^*\{L\ge16A_k\}
 \le\frac{4^{-n}}{1-2^{-n}},\qquad
 \sum_kN_{nk}\PP_{\mathrm{mix}}^*\{L_{\ell_n}\ge A_k\}\ge\frac{2^n}{8}
 \quad(n\ge3).
\]
The first sum is summable over batches. The second supplies many
opportunities for a success. Tests along the actual curve are separated
in conformal-radius time so that angle mixing gives a conditional lower
bound after each preceding test. This uses their actual stopped pasts,
not independence. One final cleanup scale in each batch pays for points
missed by all of its tests.

Writing $b_{ni}$ and $r_{ni}$ for the resulting thresholds and radii,
the gauge interpolates their costs by the quadratic expressions
$b_{ni}r_{ni}^d\max\{1,(r/r_{ni})^2\}$ and takes the infimum over
\emph{every} entry. The quadratic factor accounts for subdivision in
the plane. Radii in different batches may interleave.

For the lower bound, a strict past-mass identity uses the conditional
Green kernel to bound expected mass in overfull squares through centered
probabilities that a mass threshold is exceeded before the marked point
is reached. The upper tails make these expected losses summable over
all batch entries. Positive mass
therefore remains on a set that avoids all overfull squares after one
common batch cutoff. The resulting finite nonatomic measure is bounded
by every late quadratic expression. Its finite total mass supplies the
bounds for the finitely many excluded entries, so taking the full
infimum is legitimate. Summing this one density bound over arbitrary
countable covers proves positivity.

For the upper bound, successful finite tests force terminal-mass
candidates. A greedy selection makes their charging balls disjoint, and
one actual mass in a larger bounded region pays for their enlarged
covering balls. Only adapted test failures enter the finite change of
law; terminal candidate membership is used afterward through inclusion.
The cleanup scale covers every missed cell, including the bottom row
at the real boundary. These measurable covers have diameters tending
to zero and uniformly bounded expected cost. Their lower limiting cost
and Fatou's Lemma give \eqref{intro:box-mean}.

\subsection{Prior work and proof organization}

The Green-potential approach to trace mass has its roots in natural
parametrization. Lawler and Sheffield constructed natural parametrization
initially for $0<\kappa<4(7-\sqrt{33})$
\cite[Theorem~3.1]{LawlerSheffield2011}; Lawler and Zhou obtained the
full range $0<\kappa<8$ \cite[Theorem~1]{LawlerZhou2013}.
Lawler and Rezaei constructed $d$-dimensional Minkowski content and
identified it, up to normalization, with natural parametrization
\cite[Theorem~1.1 and following discussion]{LawlerRezaei2015}.
Our common-kernel construction proves the properties needed here
locally; identifying its mass with those measures is unnecessary.
Multipoint estimates and their higher-moment consequences also have
substantial predecessors \cite{LawlerWerness2013,RezaeiZhan2017}.
Rezaei--Zhan subsequently proved existence, local H\"older continuity,
and bounds up to constants for chordal multipoint Green functions at
every finite order \cite[Theorem~1.1]{RezaeiZhan2018Green}.
The domain-uniform bound in \eqref{intro:multipoint} also follows from
Rezaei--Zhan \cite[Theorem~1.1]{RezaeiZhan2017} and the Koebe estimates.
We give a direct product-supermartingale proof with each point frozen
after its own target hit.

The use of regions carrying large mass together with a fine residual
cover appears in Rezaei \cite[Section~3, equations~(9)--(11)]{Rezaei2018}.
Zhan \cite[Remark~4.5]{Zhan2019Optimal} gives a conditional covering
sketch for zero $d$-dimensional Hausdorff measure.
The present finite batches, actual restart identities and all-time
boundary cleanup supply the covering proof for the specified gauge.
The quadratic envelope is related to the planar regularization of
Peres--Solomyak \cite[Lemma~1.2 and Section~5]{PeresSolomyak2005};
the strict past-mass calculation uses the same remaining-potential
mechanism as Zhan's decomposition along the SLE curve
\cite[Proposition~2.2 and equations~(4.3)--(4.5)]{Zhan2019Decomposition}.

Holden and Yuan identify a positive fine-mesh variation, up to a
deterministic constant, with natural parametrization
\cite[Theorems~1.7--1.8]{HoldenYuan2026}. In the half-plane their
result includes positive-time intervals that touch the boundary.
Its ordered time partitions differ from the arbitrary spatial covers
used here. The moment prescription proves exactness for its own gauge;
it asserts neither a log-log asymptotic nor identification of Hausdorff
measure with natural parametrization.

A later companion \cite[Theorem~1.1]{OpenAI2026Explicit} gives
exactness for the regular closed-form gauge
\[
 r^d\bigl(\log\log(1/r)\bigr)^{(2-d)/2}
\]
at small radii, by quantitative mass tails and dense visits.
Its argument establishes its own stopped-law, mass and covering estimates.
The present argument leaves the asymptotic relation between that formula
and the moment-defined gauge $h_\kappa$ undetermined.

Section~\ref{op:section} proves finite centered laws, mixing and the
remaining-potential identity in the required normalization.
Sections~\ref{mp:section} and~\ref{mass:section} establish
\eqref{intro:multipoint}, the common mass kernel and actual restart.
Section~\ref{coef:section} recovers the deterministic disk moments;
Section~\ref{gauge:section} turns them into batches and the gauge.
Section~\ref{hd:section} proves the strict past-mass comparison and
arbitrary-cover lower bound. Section~\ref{hd:cover-section} constructs
the all-time upper covers and assembles the theorem.
The external geometric inputs are continuous trace generation
\cite[Theorem~5.1]{RohdeSchramm2005} and the classical Koebe estimates
\cite[Section~3.2]{Beliaev2015Notes}. Every subsequent centered-law,
multipoint, mass and covering argument is proved below.
